\section{Discussion}
\label{sec:tartisma}

The findings first give a clear answer to RQ1: the seam left by the
as-shipped Gaussian blend when truncated at a finite radius is not a
cosmetic or negligible defect but a structural one, measurable at three
orders of magnitude relative to smoothstep (approximately 3200x in C0,
530x in C1) and statistically highly significant. The source of this
defect is not a coding error but a property of radial basis functions
that has been known for more than fifty years: the Gaussian-type radial
functions defined by Hardy \citep{hardy1971multiquadric} are
infinite-support and never reach exactly zero at any finite radius;
Wendland \citep{wendland1995piecewise} proposed compact-support,
minimal-degree positive-definite radial functions as exactly the
solution to this problem. This study's contribution is to demonstrate
this fifty-year-old mathematical fact quantitatively in a concrete
engineering setting (a physically queryable lunar-terrain height field)
and to test whether it is reflected in RL training. Unlike the RL
results discussed below, this analytic finding is validated over 30
seeds and is not subject to single-run uncertainty; it is therefore the
study's most robust and least debatable finding.

In answer to RQ2, the proposed hybrid method preserves the Gaussian's
soft ``bell'' profile in the contact region and smoothstep's compact
support beyond it, while creating no new seam at its own internal
transition (switch); this is an expected result, consistent with the
partition-of-unity framework described by Ohtake et al.\
\citep{ohtake2003multiscale}: a convex combination of two C1-smooth
functions with a C1-smooth weight is itself C1-smooth. The cost of this
advantage is an approximately 4.3x higher microsecond cost than
smoothstep in computing the pure fade formula (because both branches
must be evaluated). This formula-level cost is reflected only partially
and inconsistently in the real training's wall-clock time: in the
initial (seed=42) measurement, Gaussian and hybrid were 56\% and 35\%
slower, respectively, but averaged over three seeds this gap drops to
20\% and 12.5\%, and the standard deviation (20--30\% of the mean) shows
how much a single run can deviate (\cref{sec:egitimsonuclari}). This
suggests that the real wall-clock difference stems largely from
run-to-run system variability (background load, thermal state, terrain
pool timing), with the pure formula cost contributing only a small and
consistent part; this is a separate observation, not to be
over-interpreted, but one that shows single-run measurements can also be
misleading for cost comparisons.

The answer to RQ3 changed substantially after multiple seeds, and this
change is itself one of the study's most instructive findings. In the
initial (single-seed, seed=42) measurement, Gaussian and hybrid appeared
to reach a final success rate approximately 8--12 percentage points
higher than smoothstep, and this initial draft attributed it to a
mechanism such as ``a difference in the transition-band weight
profile.'' In the measurement extended with two additional independent
seeds (seed=7, seed=123) (\cref{sec:egitimsonuclari}), the cross-method
gap shrank (the Gaussian-smoothstep gap dropped from 11\% to 4.7\%), the
hybrid-smoothstep gap changed sign, and no pairwise comparison was
statistically significant on any of the three metrics (success rate,
timeout rate, harsh-landing rate; all $p\ge 0.20$). Recalling that the
Gaussian variant used in the real training (calibrated, $k=\ln(1000)$)
already has analytic boundary continuity nearly equal to smoothstep's
(\cref{tab:continuity}), this result is in fact an expected consistency:
if there is no measurable difference in boundary continuity, it is not
surprising that there is also no strong difference in RL training
performance. The apparent gap in the initial measurement was most likely
ordinary run-to-run variance stemming from the GPU-parallel PhysX
simulation's own internal nondeterminism and the timing-dependence of
the terrain pool's background CPU generation process --- exactly the
pseudoreplication pitfall warned against by Henderson et al.\
\citep{henderson2018deep}. Because this study thereby refutes its own
initial finding, it offers a self-contained case study on how misleading
single-run RL comparisons can be: going from $n=1$ to $n=3$ was enough to
render an apparently ``significant'' finding statistically
indistinguishable. That said, $n=3$ is still low-powered, and the
present data cannot rule out the presence of a small real effect
(especially given that the Gaussian consistently has the highest
average); a larger number of seeds is needed to reach a definitive
conclusion (\cref{sec:sonuc}).

The scope of this study was deliberately kept narrow in order to isolate
the effect of the blending function as a single variable; the three
methods were compared only against each other, with the rest of the
production environment (physics engine, reward function, agent
architecture, curriculum schedule) held fixed. Bringing in another
terrain-blending method from the literature --- e.g., one implemented in
a different simulator --- as an external benchmark was deliberately not
chosen; such a comparison would make it impossible to separate the
effect of the blending function from the countless other differences in
that simulator's physics engine, observation/reward design, and agent
architecture. The cost of this choice is that the extent to which the
findings generalize beyond this specific environment (Isaac Lab/PhysX,
SAC, this task definition) --- external validity --- was not tested in
this study; this point is addressed as a concrete future-work direction
in \cref{sec:sonuc}.

The most important limitation of this study is that, with $n=3$, it
still has low statistical power: the smallest $p$-value the Mann-Whitney
test can produce is already bounded with the current data (around a
minimum of $p\approx 0.10$ for a 3-vs-3 comparison), so an ``no
significant difference'' finding does not prove that the methods are
truly equivalent --- it only shows that the seemingly large gap in the
initial measurement could not be reliably confirmed. The second
limitation is that the training scale (5.73 million steps, approximately
350 iterations) remained below the full-scale configuration (1200
iterations) owing to the computational budget; no run may have reached
full convergence, meaning the observed plateaus may reflect an
interim-stage performance difference rather than a final/asymptotic one.
The third limitation is that the analytic/CPU study isolates only the
fade formula itself and does not fully disentangle the source of the
real-training wall-clock difference (terrain generation, physics
stepping, system variability) --- the three-seed data show that this gap
shrinks, but do not conclusively identify its source.
